\section{A hierarchy of consistent block witnesses}
\label{sec:block-witness}

The pair witness can lose information because its optimal fitted values are
chosen separately for each pair. Requiring one set of fitted values to
respect every distance in a larger block removes some of this freedom. We
now make this observation precise for a \emph{scalar} target. The result
connects the pair obstruction to the optimal Lipschitz risk, while exposing
the statistical cost of using larger blocks.

Fix the representation, pseudometric, and sensitivity budget from
Section~\ref{sec:theory}. Let $(U_i,Y_i)_{i=1}^b$ be independent copies
from the full evaluation law, where $b\geq2$, and assume
$Y\in[y_-,y_+]$ almost surely. Write $D=y_+-y_-$. Define
\begin{equation}
 \begin{aligned}
 T_b&=\frac{1}{b-1}\min_{v\in\mathbb R^b}
       \sum_{i=1}^b(Y_i-v_i)^2,
 &\mathcal{W}_b&=\mathbb E T_b,\\
 &\text{subject to }|v_i-v_j|\leq\kappa d(U_i,U_j)
       \quad(i,j\leq b).
 \end{aligned}
 \label{eq:block-witness}
\end{equation}
The feasible set contains every constant vector. The minimum therefore
exists and is unique. Its value is unchanged if the fitted values are
restricted to $[y_-,y_+]$, because clipping preserves feasibility and
cannot increase squared error. The factor $1/(b-1)$ is essential: if
$a_b$ denotes the usual minimum empirical MSE, then
$T_b=b a_b/(b-1)$. This correction is justified by the freedom to translate
every fitted value by the same constant; it is not a generic correction
for arbitrary empirical-risk minima.

\begin{theorem}[Monotone block hierarchy]
\label{thm:block-hierarchy}
Under these assumptions, the following statements hold.
\begin{enumerate}[label=(\roman*),leftmargin=*,itemsep=2pt]
\item The first member is exactly the pair score:
\[
 T_2=\tfrac12\bigl(|Y_1-Y_2|-\kappa d(U_1,U_2)\bigr)_+^2.
\]
Consequently $\mathcal{W}_2=W_{\kappa,A}$ when $A$ is constant.
\item For every $b\geq2$,
\begin{equation}
 0\leq T_b\leq B_b^{\rm blk},
 \qquad B_b^{\rm blk}=\frac{D^2\lfloor b^2/4\rfloor}{b(b-1)},
 \qquad \mathcal{W}_b\leq R^*_\kappa.
 \label{eq:block-range}
\end{equation}
\item If $2\leq m\leq b$ and $T_m(S)$ is the score computed from a
subset $S\subseteq\{1,\ldots,b\}$ of size $m$, then
\begin{equation}
 \binom bm^{-1}\sum_{|S|=m}T_m(S)\leq T_b.
 \label{eq:block-subset}
\end{equation}
In particular,
$\mathcal{W}_2\leq \mathcal{W}_3\leq\cdots\leq R^*_\kappa$.
\item If the representation takes values in a compact metric space
and $d$ is that metric, then $\mathcal{W}_b\uparrow R^*_\kappa$ as
$b\to\infty$.
\end{enumerate}
\end{theorem}

\begin{proof}
For two observations, let
$\Delta=|Y_1-Y_2|$ and $c=\kappa d(U_1,U_2)$. If $\Delta\leq c$,
both targets can be fitted exactly. Otherwise the optimal fitted
separation is $c$, with the two residuals having magnitude
$(\Delta-c)/2$ and opposite signs. Their squared sum is
$(\Delta-c)^2/2$, proving (i).

For any fixed admissible decoder $g$ of finite risk, put
$r_i=Y_i-g(U_i)$ and $\bar r=b^{-1}\sum_i r_i$.
The translated values $v_i=g(U_i)+\bar r$ are feasible. Thus
\[
 T_b\leq\frac1{b-1}\sum_i(r_i-\bar r)^2,
 \qquad
 \mathcal{W}_b\leq\operatorname{Var}(Y-g(U))\leq R(g\circ\Gamma).
\]
Taking the infimum over $g$ proves the population bound. Using the
constant fitted value $\bar Y$ instead bounds $T_b$ by the unbiased
sample variance of the targets. This variance is a convex function
of each target coordinate, so its maximum over $[y_-,y_+]^b$ occurs
at interval endpoints. If $k$ coordinates equal $y_+$ and $b-k$
equal $y_-$, its value is $D^2 k(b-k)/\{b(b-1)\}$. Maximizing over
integer $k$ proves the range bound.

For (iii), let $v$ minimize the full-block objective and set
$r_i=Y_i-v_i$. Translation invariance implies $\sum_i r_i=0$:
otherwise translating $v$ by the residual mean would strictly lower
the objective. Hence
\begin{equation}
 T_b=\frac1{b(b-1)}\sum_{i<j}(r_i-r_j)^2.
 \label{eq:block-residual-identity}
\end{equation}
On a subset $S$, restricting $v$ and translating by that subset's
residual mean gives feasible fitted values. Therefore
\[
 T_m(S)\leq\frac1{m(m-1)}
                  \sum_{i<j,\ i,j\in S}(r_i-r_j)^2.
\]
A uniformly chosen $m$-subset contains a given pair with probability
$m(m-1)/\{b(b-1)\}$. Averaging proves
\eqref{eq:block-subset}; taking expectations proves monotonicity.
This is an average-over-subsets statement, not an ordering between
arbitrary realized blocks of different sizes.

To prove (iv), any feasible vector extends to a scalar Lipschitz
function by
$g_v(u)=\min_i\{v_i+\kappa d(u,U_i)\}$.
The triangle inequality and the finite constraints show that this
function is $\kappa$-Lipschitz and equals $v_i$ at $U_i$.
Clipping it to $[y_-,y_+]$ preserves those properties. Thus $a_b$
is the empirical minimum over the same bounded Lipschitz class as
the population problem. On a compact metric space this class is
compact in the uniform norm: it is uniformly bounded and
equicontinuous. For each fixed $\epsilon>0$, a finite uniform
$\epsilon$-net and the ordinary strong law give uniform convergence
of empirical risks over that net. Since
$|(y-f(u))^2-(y-g(u))^2|\leq2D\|f-g\|_\infty$ for interval-valued
$f,g$, the convergence extends to the entire class. Therefore
$a_b\to R^*_\kappa$ almost surely. Multiplication by $b/(b-1)$
and the bound $T_b\leq D^2/2$ give $\mathcal{W}_b\to R^*_\kappa$ by
dominated convergence.
\end{proof}

The result strengthens the population witness without changing the decoder
class. In particular, every block score dominates the average of its
pair scores. The additional information is joint consistency of the fitted
values across pairs, rather than a new concentration inequality.
Optimization of empirical $U$-statistics is established
\citep{clemencon2016scaling}: equation~\eqref{eq:block-residual-identity}
places this construction in that framework. Lower bounds obtained by
averaging empirical optimization values over independent batches are also
established in stochastic programming
\citep{linderoth2006empirical,guigues2017confidence}. The specific
connection here is the centered normalization, its exact pair-obstruction
member, and the monotone hierarchy for the Lipschitz prediction problem.

\subsection{Exact gaps and the limits of fixed block sizes}

Consistency for a fixed distribution does not provide a universal choice
of $b$. Two exact calculations distinguish rare predictive conflicts from
the capacity to interpolate many distinct inputs.

\begin{proposition}[Two exact block-size laws]
\label{prop:block-rare}
For the three-atom distribution from Section~\ref{sec:theory}, with
$(U,Y)=(3,0),(0,-1),(1,1)$ of probabilities
$1-2\varepsilon,\varepsilon,\varepsilon$ and $\kappa=1$,
\begin{equation}
 R^*_\kappa=\frac\varepsilon2,
 \qquad
 \mathcal{W}_b=\frac{2b\varepsilon-1+(1-2\varepsilon)^b}{4(b-1)}.
 \label{eq:block-rare-exact}
\end{equation}
For every fixed $b$, $\mathcal{W}_b/R^*_\kappa\to0$ as
$\varepsilon\downarrow0$.

Separately, let $U$ be uniform on $M\geq1$ points whose pairwise
distances are at least $D/\kappa$, where $D,\kappa>0$, and let
$Y\in\{0,D\}$ be independent of $U$ and equiprobable. Then
\begin{equation}
 R^*_\kappa=\frac{D^2}{4},
 \qquad
 \mathcal{W}_b=\frac{D^2}{4(b-1)}
       \left[b-M\{1-(1-1/M)^b\}\right].
 \label{eq:block-occupancy}
\end{equation}
For fixed $b$, $\mathcal{W}_b\to0$ as $M\to\infty$, while the optimal risk
stays $D^2/4$.
\end{proposition}

\begin{proof}
In the first law, let $N_-$ and $N_+$ count the two rare atoms in
a block, and put $N=N_-+N_+$. The common atom can be fitted at zero
without restricting the optimum at the rare atoms. Their constrained
weighted least-squares error is $N_-N_+/N$ when $N>0$, and zero
otherwise. Conditional on $N$, $N_-$ is binomial with parameters
$N$ and $1/2$, so the conditional mean of this error is
$(N-1)/4$ for $N\geq1$. Since
$N\sim\operatorname{Binomial}(b,2\varepsilon)$,
\[
 (b-1)\mathcal{W}_b=\tfrac14\{\mathbb EN-\Pr(N\geq1)\}
 =\tfrac14\{2b\varepsilon-1+(1-2\varepsilon)^b\}.
\]
Its expansion at zero gives
$\mathcal{W}_b/R^*_\kappa=b\varepsilon+O(\varepsilon^2)$.

In the second law, any fitted values in $[0,D]$ at distinct support
points satisfy the Lipschitz constraints. At a support point observed
$n>0$ times, the empirical optimum is its target mean, with expected
squared residual sum $(n-1)D^2/4$. If $O_b$ is the number of occupied
support points, the conditional expected total is
$(b-O_b)D^2/4$. Each of the $M$ points is occupied with probability
$1-(1-1/M)^b$. Taking expectations gives
\eqref{eq:block-occupancy}. Independence of $Y$ and $U$ gives the
population optimum $D^2/4$, attained by the constant $D/2$.
\end{proof}

The second example consists entirely of finite, hence compact, metric
spaces. It rules out a distribution-free additive approximation rate
depending only on $b$ and the target range over all such spaces. The
support points can be $u_j=D e_j/(\kappa\sqrt2)$ in $\mathbb R^M$,
where $e_j$ are the coordinate unit vectors. Their diameter is
$D/\kappa$ for every $M\geq2$, so a bound on representation diameter
alone does not remove this limitation. For fixed $b$,
$\mathcal{W}_b/R^*_\kappa=b/(2M)+O(M^{-2})$ as $M\to\infty$.
The
first example shows why increasing block size can be necessary even on
three fixed support points. Larger blocks address these losses, but do
not make an uninformative certificate evidence of adequate representation.

\subsection{Calibration and certified optimization values}

Use disjoint blocks of $b$ independent calibration units, giving
$K_b=\lfloor n/b\rfloor$ independent scores. This makes the effect of
block size explicit: increasing $b$ tightens $\mathcal{W}_b$ but decreases the
number of independently calibrated scores. The hierarchy therefore
does not order finite-sample confidence endpoints.

The quadratic program in \eqref{eq:block-witness} must be solved from
below for a lower certificate. Write its inequalities as $Av\leq c$,
where rows of $A$ are coordinate differences and
$c_{ij}=\kappa d(U_i,U_j)$. For every nonnegative multiplier vector
$\lambda$, weak duality gives the lower value
\begin{equation}
 d_b(\lambda)=\lambda^\top(AY-c)
                       -\tfrac14\|A^\top\lambda\|_2^2
 \ \leq\ \min_{Av\leq c}\|Y-v\|_2^2.
 \label{eq:block-dual}
\end{equation}
This follows by minimizing
$\|Y-v\|_2^2+\lambda^\top(Av-c)$ over $v$; its minimizer is
$Y-A^\top\lambda/2$.
Thus
$\underline T_b=\max\{0,d_b(\lambda)/(b-1)\}\leq T_b$.
A primal value alone has the wrong inequality direction. The
implementation reports primal feasibility and the primal--dual gap;
floating-point tolerances are numerical checks rather than proofs by
verified arithmetic.

\begin{corollary}[Independent-block lower calibration]
\label{cor:block-calibration}
Suppose the same fixed measurable procedure computes each score
$\underline T_{b,k}\in[0,B_b^{\rm blk}]$ from block $k$, with
$\underline T_{b,k}\leq T_{b,k}$, and $K_b\geq1$.
For an allocated failure probability $\alpha\in(0,1)$ and $B_b^{\rm blk}>0$,
let
\begin{equation}
 L_b=B_b^{\rm blk}\operatorname{kl}_-\!\left(
       \frac{K_b^{-1}\sum_{k=1}^{K_b}\underline T_{b,k}}{B_b^{\rm blk}},
       \frac{\log(1/\alpha)}{K_b}\right),
 \label{eq:block-lower}
\end{equation}
where $\operatorname{kl}_-$ is defined in
Appendix~\ref{app:kl-calibration}. Then
$L_b\leq \mathcal{W}_b\leq R^*_\kappa$ with probability at least
$1-\alpha$. If $B_b^{\rm blk}=0$, take $L_b=0$. A finite family of block
sizes, representations, horizons, and comparison methods can be
selected after calibration when the sum of their allocated failure
probabilities is at most the declared total.
\end{corollary}

\begin{proof}
The bounded-mean inequality \eqref{eq:kl-bounded-mean} gives
$L_b\leq\mathbb E\underline T_{b,1}$ with failure probability
at most $\alpha$. The pointwise score inequality and
Theorem~\ref{thm:block-hierarchy} give
$\mathbb E\underline T_{b,1}\leq \mathcal{W}_b\leq R^*_\kappa$.
A union bound proves the finite-family assertion; different methods
may reuse the same calibration observations.
\end{proof}

The empirical Bernstein radius can be substituted when $K_b\geq2$.
Block sizes and optimization rules are fixed before calibration; taking
the largest endpoint across a frozen family requires the stated joint
allocation. This is distinct from multiplying the full-sample
empirical-minimum endpoint in Appendix~\ref{app:erm-comparator} by
$n/(n-1)$, which that endpoint's proof does not justify.

\paragraph{Conditioning and forecasting origins.}
For a common origin or stratum $A$, draw the $b$ observations independently
conditional on $A$. The same proof gives a monotone hierarchy
$\mathcal{W}_{b,A}\leq R^*_\kappa$. For finitely many values of $A$ and the
compact-support conditions above, its limit is
\[
 \sum_a\Pr(A=a)\inf_{g\in\mathcal G_\kappa(d)}
                 \mathbb E[(Y-g(U))^2\mid A=a],
\]
which can be strictly below the optimum of one shared decoder. The
decoder attaining each conditional infimum may differ with $a$.
Consequently, averaging a block score over common forecast origins
preserves validity and monotonicity but does not generally establish
convergence to the shared-decoder risk. For that consistency claim,
one may instead draw an independent uniform origin in each independent
episode; the resulting observations are independent draws from the full
evaluation law. Neither construction treats overlapping windows as
independent calibration units.
